\documentclass[12pt]{article}
\usepackage[utf8]{inputenc}
\usepackage[T1]{fontenc}
\usepackage{hyperref}
\title{Punic Language Resources for K-12 Curriculum Development}
\author{Compiled by Claude Sonnet 5.5 (claude-sonnet-5-5) for Andr\'as Kornai}
\date{}
\begin{document}
\maketitle

\section{Introduction}
Punic (ISO 639-3 xpu; Glottocode puni1241) is the extinct Phoenician language of Carthage and the western Mediterranean. It belongs to the Afro-Asiatic family. It is documented through inscriptions, including Neo-Punic texts from Roman Tripolitania, and through fragments of Punic speech in Plautus's \emph{Poenulus}. The search found only scholarly resources. It found no school programme.

\section{Primary Educational Resources}
No primers or school programmes for Punic were found. The best candidate for adaptation is the grammar below.

\section{Linguistic and Reference Materials}
\begin{itemize}
\item Charles Krahmalkov, \emph{A Phoenician-Punic Grammar}, Brill, 2001 (SBL Press reprint, 2014). It covers phonology, orthography, morphology and syntax, and includes the Poenulus passages. Catalogue records: \url{https://library.arce.org/portal/Record/40648}, with its table of contents at \url{https://library.arce.org/portal/Record/40648/TOC}. Publisher pages: \url{https://www.degruyterbrill.com/document/isbn/9789004294202/html} and \url{https://cart.sbl-site.org/books/069556P}.
\item Donner and R\"ollig, \emph{Kanaan\"aische und Aram\"aische Inschriften} (KAI), a standard inscription sourcebook mentioned in a university course syllabus. No URL for the book itself was found.
\item J. C. L. Gibson, \emph{Textbook of Syrian Semitic Inscriptions}, vol. 3 (Phoenician). It was mentioned in the search results, with no URL.
\end{itemize}

\section{Religious and Community Resources}
None found.

\section{Digital and Audio Resources}
The search results also listed a general Academia.edu topic page: \url{https://www.academia.edu/Documents/in/Punic_Language}. No audio or software resources were found.

\section{Curriculum Development Context}
A school unit could teach the Phoenician-Punic alphabet and simple inscriptions. Krahmalkov's grammar would serve as the teacher's reference. A full immersive curriculum is not realistic.
\end{document}